\chapter{Implementation}
\label{ch:implementacao}

\section{Introduction}

This chapter describes the implemented system and explains the reasoning behind major design
choices. For each component, two perspectives are presented: \textbf{functional behaviour} (what
the system does) and \textbf{design rationale} (why that approach was adopted and which
trade-offs it entails). Architectural context and diagrams appear in
Chapters~\ref{ch:engenharia} and~\ref{ch:estado-arte}; validation interpretation appears in
Section~\ref{sec:discussion-results}.

The primary use case---an analyst submitting a suspicious Windows executable---proceeds as follows:
the system validates the upload, assigns a unique job identifier, runs static analysis (with optional
real-time streaming on the web), optionally schedules sandbox execution, and presents unified results
with decompiled code, behavioural findings, and an explainable risk score. The analyst may
optionally query the Gemini assistant about suspicious pseudo-C excerpts.

\section{Backend - Static Analysis Pipeline}

\textbf{Functional behaviour.}
The static analysis engine orchestrates seven phases: (1)~PE metadata extraction;
(2)~managed decompilation via ILSpy when the sample is a .NET assembly; (3)~import, string, and
evasion heuristics; (4)~deobfuscation of common encodings; (5)~YARA rule matching;
(6)~aggregation into a 0--100 risk score with five severity levels; and (7)~structured report
generation, including native disassembly and Ghidra pseudo-C when ILSpy cannot produce managed
code. Phases run sequentially on the host; results are persisted per job and exposed through a
REST API that supports synchronous analysis, NDJSON streaming for live progress, asynchronous
job creation for combined static/dynamic modes, artefact retrieval, and storage management. Web
uploads of \texttt{.cs} source files receive string heuristics only; full PE analysis requires a
compiled binary.

\textbf{Design rationale.}
Phases are ordered \emph{cheap-before-expensive} because YARA compilation and Ghidra dominate
wall-clock time on typical laboratory hardware (Table~\ref{tab:benchmark}). Running lightweight
checks first lets the pipeline fail fast on obvious benign samples and surface early indicators
before committing to decompilation---a pattern aligned with professional triage workflows
\cite{bensaoud2024,anderson2012}. Routing .NET samples through ILSpy before Ghidra reflects the
prevalence of managed RAT loaders and avoids unnecessary native analysis when managed IL is
available. Limiting \texttt{.cs} handling to heuristics is a deliberate safety trade-off: compiling
untrusted source on the server would introduce build-time risk without matching the PE analysis
depth students need after compilation.

The API offers both synchronous and job-based modes because use cases differ. Quick CLI or API
calls benefit from a single request--response; the web interface benefits from streaming logs
during long Ghidra runs so users perceive progress rather than a frozen screen. Job identifiers
unify static and dynamic artefacts when both modes run on the same sample---essential for the
combined \texttt{both} workflow. Token authentication on analysis routes follows the fail-closed
policy defined in Chapter~\ref{ch:engenharia}; OpenAPI metadata supports integration without
requiring this report to catalogue every route.

\section{Web Frontend}

\textbf{Functional behaviour.}
The web application provides drag-and-drop upload, mode selection (\texttt{static},
\texttt{dynamic}, or \texttt{both}), real-time progress during static analysis, job submission for
Path~A sandbox runs, and tri-panel visualisation: structured report, Ghidra pseudo-C, and ILSpy IL
for .NET samples. Cross-reference navigation links indicators to decompiled regions; bilingual
support (Portuguese/English) covers classroom use at \ubi{}.

\textbf{Design rationale.}
A browser-first interface was chosen so students need no local installation beyond the backend and
sandbox: any laboratory PC with a modern browser can submit samples and read explainable reports.
Streaming progress addresses a recurring usability complaint in long-running analysis tools---users
abandon workflows they cannot monitor. Splitting results into three panels reflects how analysts
actually work: triage from the report and score, then drill into code for confirmation. Keeping
static visualisation in the browser while WPF handles Path~B avoids duplicating the rich results UI
in two codebases; both entry points converge on the same job store and permalink model.

Figure~\ref{fig:web-upload} shows the \emph{Drop \& Analyze} upload zone with analysis mode
selection and real-time progress indicators.

\begin{figure}[H]
    \centering
    \reportfig{fig-5-1-web-upload.png}
    \caption{Web interface - upload zone (\emph{Drop \& Analyze})}
    \label{fig:web-upload}
\end{figure}

Figure~\ref{fig:web-results} presents the results view: risk score, structured report, and
side-by-side pseudo-C and IL panels for decompiled .NET code.

\begin{figure}[H]
    \centering
    \reportfig{fig-5-2-web-results.png}
    \caption{Web interface - report, pseudo-C/IL, and risk score}
    \label{fig:web-results}
\end{figure}

\subsection{Gemini Assistant in the Pseudo-C Panel}

\textbf{Functional behaviour.}
In the pseudo-C panel, the user may open a chat dialog that sends the \emph{visible} decompiled
excerpt (bounded by a character limit) to Google's Gemini API, together with a natural-language
question. The API key is stored in browser \texttt{localStorage}; a demonstration mode returns
simulated answers without external calls.

\textbf{Design rationale.}
The assistant was placed client-side so sample excerpts and third-party credentials never pass
through the FastAPI backend---a confidentiality decision consistent with local academic use
\cite{galli2024}. This sacrifices centralised audit logging and model governance in exchange for
user control over what code leaves the machine. Bounding the excerpt limits accidental exfiltration
of entire binaries through the chat surface. The demo mode supports classroom presentations on
networks where external API access is blocked.

\section{WPF Desktop Application}

\textbf{Functional behaviour.}
The WPF client (\texttt{Rat Analyzer.exe}) bootstraps backend and frontend dependencies, persists
language preference, and exposes a dashboard for drag-and-drop analysis. Static analysis opens the
web results view in the system browser; behavioural analysis invokes Path~B PowerShell automation
and, when a prior static job exists, publishes the sandbox report to the same job identifier so
static and dynamic findings appear together.

\textbf{Design rationale.}
WPF was retained as a desktop entry point because Hyper-V orchestration on Windows often requires
elevated privileges and direct PowerShell integration that a pure web page cannot safely perform.
Delegating static visualisation to the browser reuses the full React results experience instead of
reimplementing panels in XAML---a pragmatic split that trades a two-window workflow for faster
development and a single source of truth for report rendering. Publishing Path~B results back to
the shared job store preserves the permalink model students already use from the web UI.

Figure~\ref{fig:wpf-dashboard} shows the \texttt{MainDashboardView}: drag-and-drop sample entry,
static versus behavioural analysis selection, and integration with Path~B sandbox orchestration.

\begin{figure}[H]
    \centering
    \reportfig{fig-5-3-wpf-dashboard.png}
    \caption{WPF desktop application - main screen}
    \label{fig:wpf-dashboard}
\end{figure}

\section{VM Agent and Hyper-V Sandbox}

Dynamic analysis is documented in Figures~\ref{fig:sandbox-paths}--\ref{fig:seq-path-b} and
Section~\ref{sec:caminhos-ab}. Two orchestration paths coexist; each solves a different integration
problem.

\subsection{Path A - Backend + HTTP vm-agent}

\textbf{Functional behaviour.}
With the Hyper-V driver enabled, the Python backend restores the \texttt{CleanState} snapshot,
starts the guest VM, uploads the sample to an in-guest HTTP agent, triggers execution, polls for
completion, and retrieves a JSON behavioural report. The agent authenticates requests with a
shared token, confines uploads to a safe directory, caps file size, and serialises runs so only one
sample executes at a time.

\textbf{Design rationale.}
Path~A integrates dynamic analysis into the same job queue as static analysis---the natural choice
for the web \texttt{dynamic}/\texttt{both} modes. An HTTP agent inside the guest decouples Python
orchestration from guest OS details: the host only needs snapshot management and REST calls, which
simplifies testing with a stub driver during development. Single-run serialisation prioritises
isolation over throughput in teaching labs where concurrent malware execution adds little pedagogical
value but increases cross-contamination risk \cite{alrawi2024}. Lighter telemetry than Path~B is an
accepted limitation: Path~A optimises for API cohesion, not maximum behavioural depth.

\subsection{Path B - PowerShell Scripts}

\textbf{Functional behaviour.}
Path~B is driven by the WPF client and a PowerShell pipeline (phases A--G): VM setup preflight,
sample copy into the guest, execution of a guest-side collector (\texttt{Run-MalwareAnalysis.ps1}),
diff of files, registry, network, and Sysmon events when available \cite{sysmon2022}, then report
transfer via PowerShell Direct or \texttt{Copy-VMFile}. The guest writes a completion digest with
report SHA-256 and byte length; the host accepts the copy only if both match
(Listing~\ref{lst:sha256-verify}).

\textbf{Design rationale.}
Path~B exists because rich Windows-native instrumentation predated the HTTP agent and because
PowerShell Direct is the most direct way to reach a Hyper-V guest from an elevated desktop session.
SHA-256 verification across the VM boundary addresses a real failure mode: truncated or partial
file copies that would otherwise produce silent false negatives in behavioural reporting. Keeping
Path~B outside the Python driver lets the sandbox evolve---new Sysmon rules, longer timeouts,
extra collectors---without redeploying the backend. The cost is duplicated orchestration logic and
two maintenance surfaces, judged acceptable for an academic platform that values depth on the
teaching path (WPF) and breadth on the API path (web).

\begin{lstlisting}[style=powershell,caption={[Post-transfer SHA-256 verification]{Post-transfer SHA-256 verification --- \filepath{SandboxCommon/VmFileTransfer.ps1}}},label={lst:sha256-verify}]
function Assert-SandboxCopiedReportHash {
    param([string] $HostReportPath, [string] $ExpectedSha256, [long] $ExpectedBytes = 0)
    if ([string]::IsNullOrWhiteSpace($ExpectedSha256)) {
        throw "Guest SHA-256 digest missing; cannot validate copied report."
    }
    if ($ExpectedBytes -gt 0) {
        $hostBytes = [long](Get-Item -LiteralPath $HostReportPath).Length
        if ($hostBytes -ne $ExpectedBytes) { throw "Copied report size does not match." }
    }
    return (Assert-FileSha256 -Path $HostReportPath -ExpectedSha256 $ExpectedSha256 ...)
}
\end{lstlisting}

Figure~\ref{fig:hyperv-vm} shows the \texttt{MalwareSandbox} virtual machine and the
\texttt{CleanState} snapshot used for isolated sample execution.

\begin{figure}[H]
    \centering
    \reportfig{fig-5-4-hyperv-vm.png}
    \caption{Hyper-V - VM \texttt{MalwareSandbox} and snapshot \texttt{CleanState}}
    \label{fig:hyperv-vm}
\end{figure}

\section{User Manual}

Operational setup (dependency installation, backend and frontend startup, WPF administrator
privileges for Hyper-V, one-time sandbox provisioning, and CLI static analysis) is documented in
the repository README and sandbox setup guide. This report omits step-by-step command listings to
keep focus on design and evaluation; students should follow versioned documentation bundled with
the source release.

\section{Experimental Evaluation}

This section reports how \rat\ was validated under the safety constraints of an academic
laboratory. The subsections below cover methodology and ethical boundaries, static scoring
boundary cases, proxy detection metrics on synthetic profiles, interpretation of the results,
benign-corpus specificity, public IOC metadata, dynamic sandbox smoke tests, and formal
performance benchmarks. Together they bound what can and cannot be claimed from the current
evaluation evidence.

\subsection{Methodology and Ethical Considerations}

Because executing real malware on university infrastructure introduces safety, legal, and
operational concerns---potential lateral movement if isolation fails, accidental distribution of
live specimens, and policy restrictions on storing malicious binaries---validation deliberately
avoids hosting active RAT executables on the development workstation. This constraint is common in
academic malware engineering: the goal is to demonstrate \emph{analytical correctness} and
\emph{pipeline integrity} without turning the laboratory into a malware repository.

The evaluation strategy therefore combines four complementary strands: (i)~\textbf{220 automated
tests} exercising API contracts, scoring logic, sandbox drivers, and UI components; (ii)~\textbf{59
stratified synthetic profiles} that encode representative benign and malicious indicator
combinations aligned with the scoring implementation; (iii)~\textbf{eight compiled benign PEs}
measuring false-positive behaviour on realistic harmless programs; and (iv)~\textbf{public IOC
metadata} from MalwareBazaar cross-referenced without downloading binaries, plus dynamic execution
of a dedicated benign smoke-test binary on Path~B. Synthetic validation is appropriate here
because the research contribution centres on \emph{explainable aggregation} of indicators: if
controlled inputs with known ground-truth labels are classified consistently, the scoring model
behaves as specified; generalisation to unseen families is a separate empirical question requiring
curated malware corpora under institutional approval.

An optional MalwareBazaar static batch harness was implemented but not executed on the host for
the same safety reasons; it remains available for supervised laboratory runs. No claim is made
regarding real-world detection rate or prevalence-weighted accuracy.

\subsection{Static Scoring Validation}

Table~\ref{tab:avaliacao-scoring} summarises representative scoring scenarios exercised
automatically in CI. These cases probe boundary behaviour: benign validation hashes must not reach
HIGH without strong evidence; isolated heuristic noise must not trigger high severity; synthetic
RAT profiles with C2, persistence, and YARA matches must classify as HIGH or CRITICAL.

\begin{table}[htbp]
    \centering
    \caption{Static scoring validation scenarios}
    \label{tab:avaliacao-scoring}
    \footnotesize
    \setlength{\tabcolsep}{4pt}
    \begin{tabularx}{\linewidth}{@{}%
        >{\raggedright\arraybackslash}p{3.2cm}%
        >{\raggedright\arraybackslash}p{2.1cm}%
        >{\centering\arraybackslash}p{0.9cm}%
        >{\centering\arraybackslash}p{1.8cm}%
        Y@{}}
        \toprule
        \textbf{Scenario} & \textbf{Source} & \textbf{Score} & \textbf{Level} & \textbf{Notes} \\
        \midrule
        \filepath{BenignVmTest.exe} &
        Benign binary &
        $<48$ &
        LOW/\allowbreak MEDIUM &
        Known hash; \texttt{validation\_sample} annotation; no YARA matches \\
        Heuristic noise &
        Synthetic profile (pytest) &
        $<48$ &
        $\neq$ HIGH &
        Isolated evasion/noise does not raise to HIGH without strong signals \\
        Synthetic RAT profile &
        Synthetic profile (pytest) &
        $\geq48$ &
        HIGH/\allowbreak CRITICAL &
        C2 imports, persistence, YARA \texttt{RAT\_Generic} \\
        YARA rules v2.0.0 &
        CI (\texttt{test\_yara\_rules}) &
        --- &
        --- &
        Four files compile without error \\
        \bottomrule
    \end{tabularx}
\end{table}

\subsection{Proxy Detection Metrics}

Proxy precision and recall were computed on \textbf{59 synthetic} profiles stratified across five
severity levels. A positive triage prediction is defined as HIGH or CRITICAL, mirroring operational
use where only elevated scores warrant sandbox execution or manual deep dive.

\begin{table}[htbp]
    \centering
    \caption{Proxy detection metrics on labelled synthetic profiles ($n=59$)}
    \label{tab:detection-metrics}
    \small
    \begin{tabular}{@{}>{\raggedright\arraybackslash}p{3.2cm}rrrrrr@{}}
        \toprule
        \textbf{Metric} & \textbf{TP} & \textbf{TN} & \textbf{FP} & \textbf{FN} &
        \textbf{Prec.} & \textbf{Rec./F1} \\
        \midrule
        HIGH/CRITICAL vs.\ below & 22 & 37 & 0 & 0 & 1.00 & 1.00 / 1.00 \\
        \bottomrule
    \end{tabular}
\end{table}

\noindent\small
Accuracy = 1.00; level accuracy (expected vs.\ predicted level) = 59/59 = 1.00.

The profile set spans 19 VERY LOW, 6 LOW, 12 MEDIUM, 14 HIGH, and 8 CRITICAL cases---from empty
profiles and Microsoft URLs through AsyncRAT/QuasarRAT-style indicator bundles with packers and
evasion. All 59 predicted levels matched expectations (37 benign negatives, 22 malicious positives).

\subsection{Discussion of Evaluation Results}
\label{sec:discussion-results}

Perfect accuracy (1.00) on synthetic profiles does \textbf{not} imply perfect malware detection in
the field. The metric measures \emph{internal consistency}: given JSON indicator bundles constructed
to exercise the same scoring functions used in production, the classifier reproduces the intended
label and severity level. Several sources of optimistic bias must be acknowledged.

First, \textbf{construct bias}: profiles were authored alongside the scorer implementation and
pytest suite; they encode the designer's mental model of HIGH/CRITICAL gates (\emph{strong\_evidence}
for YARA or combined C2/function signals). Real malware may present partial, obfuscated, or novel
combinations that fall between designed buckets. Second, \textbf{spectrum bias}: the benign side
includes deliberately weak noise profiles but not adversarially crafted binaries designed to evade
heuristics while remaining malicious \cite{imran2024}. Third, \textbf{absence of prevalence}: metrics are unweighted
across severity levels; operational false-positive cost on rare HIGH alerts differs from bulk
benign traffic in enterprise settings.

The benign corpus (specificity 1.00, zero false positives on eight PEs) supports a narrower claim:
common harmless patterns---HTTP to microsoft.com, registry reads, WinForms UI, crypto digests---do
not spuriously trigger high severity. Eight samples cannot bound specificity for all legitimate
software; network libraries and remote-administration tools remain plausible false-positive sources,
as noted in Chapter~\ref{ch:conclusoes}.

Why accuracy reaches 1.00 here is therefore straightforward: labels are defined on the same feature
schema the scorer consumes, and gates were tuned until pytest and profile suites agreed. The
scientifically honest interpretation is \textbf{verification of specification}, not \textbf{validation
of discovery}. Demonstrating the latter would require executing or statically analysing withheld
malware families under ethical approval, comparing against VirusTotal or sandbox ground truth, and
reporting confidence intervals---left as future work. Nonetheless, specification verification is
valuable: it proves the explainable pipeline behaves deterministically, a prerequisite before any
field deployment or student-facing grading.

Proxy metrics complement the benign corpus and 220 automated CI tests; together they bound
\textbf{false-positive risk on harmless code} and \textbf{logical correctness on labelled
constructs}, while explicitly not bounding \textbf{false-negative risk on live RAT campaigns}.

\subsection{Benign Corpus Static Evaluation}

Eight harmless PEs covering console I/O, file access, HTTP, registry reads, cryptography, environment
inspection, GUI, and sandbox smoke behaviour were analysed statically. All scored below the HIGH
threshold (specificity 1.00, zero false positives). Table~\ref{tab:benign-corpus} lists individual
scores; most cluster at 9 (VERY LOW), with the intentional validation sample at 16 (LOW).

\begin{table}[htbp]
    \centering
    \caption[Benign corpus static evaluation ($n=8$ PEs)]{Benign corpus static evaluation ($n=8$ PEs; measured on development host)}
    \label{tab:benign-corpus}
    \footnotesize
    \setlength{\tabcolsep}{4pt}
    \begin{tabularx}{\linewidth}{@{}%
        >{\raggedright\arraybackslash}p{3.4cm}%
        >{\raggedright\arraybackslash}p{2.1cm}%
        >{\centering\arraybackslash}p{0.9cm}%
        >{\centering\arraybackslash}p{1.6cm}%
        Y@{}}
        \toprule
        \textbf{Executable} & \textbf{Subfolder} & \textbf{Score} & \textbf{Level} & \textbf{Behaviour} \\
        \midrule
        \filepath{BenignHello.exe} & \filepath{hello/} & 9 & VERY LOW & Minimal console \\
        \filepath{BenignFileIo.exe} & \filepath{file-io/} & 9 & VERY LOW & Read/write in TEMP \\
        \filepath{BenignHttpMs.exe} & \filepath{http-ms/} & 9 & VERY LOW & HTTP GET microsoft.com \\
        \filepath{BenignRegistryRead.exe} & \filepath{registry-read/} & 9 & VERY LOW & HKCU read (no Run key) \\
        \filepath{BenignCryptoDigest.exe} & \filepath{crypto-digest/} & 9 & VERY LOW & Local SHA-256 \\
        \filepath{BenignEnvDump.exe} & \filepath{env-dump/} & 9 & VERY LOW & Environment variables \\
        \filepath{BenignGuiMsg.exe} & \filepath{gui-msg/} & 9 & VERY LOW & WinForms MessageBox \\
        \filepath{BenignVmTest.exe} & \filepath{benign-vm-test} & 16 & LOW & Sandbox smoke test (\texttt{validation\_sample}) \\
        \bottomrule
    \end{tabularx}
\end{table}

\noindent\small
True negatives = 8; false positives = 0; specificity = 1.00; all passed.

\subsection{Public IOC Metadata Reference}

Public MalwareBazaar listings for QuasarRAT, AsyncRAT, and the benign validation hash provide
external family labels without storing binaries locally. Validation uses synthetic profiles that
mirror typical indicators from those families rather than executing the listed samples.

\subsection{Dynamic Sandbox Validation}

Path~B dynamic validation confirmed snapshot restore, guest execution, completion digest emission,
PsDirect report copy, and SHA-256 verification using a dedicated benign smoke-test binary---exercising
the full isolation loop without malicious payload. Path~A confirmed execution serialisation and exit
code reporting; richer in-guest telemetry remains future work (Section~\ref{ch:conclusoes}).

\subsection{Formal Performance Benchmark}

Table~\ref{tab:benchmark} presents a reproducible benchmark on the development workstation
(Intel Core i7, 16~GB RAM, Windows~10). Static runs used five iterations; the first run's longer
Ghidra cold-cache probe explains the maximum outlier (68.4~s vs.\ 36.1~s median). Path~B completed
faster than Path~A in median terms despite richer telemetry, reflecting shorter default execution
timeouts and different observation depth---a reminder that wall-clock time trades off against
behavioural coverage.

\begin{table}[htbp]
    \centering
    \caption{Formal benchmark - static pipeline and dynamic sandbox paths}
    \label{tab:benchmark}
    \small
    \begin{tabularx}{\linewidth}{@{}>{\raggedright\arraybackslash}p{3.2cm}ccp{1.0cm}Y@{}}
        \toprule
        \textbf{Operation} & \textbf{$n$} & \textbf{Median} & \textbf{IQR} & \textbf{Notes} \\
        \midrule
        Static pipeline (full) & 5 & 36.1~s & 3.0~s & Min 33.1~s, max 68.4~s; first run includes cold-cache Ghidra probe \\
        VM restore + boot & 3 & 52~s & 8~s & \texttt{Restore-VMSnapshot} + \texttt{Start-VM} \\
        Path A (HTTP vm-agent) & 3 & 198~s & 22~s & Includes 120~s sample timeout cap; basic telemetry \\
        Path B (PsDirect) & 3 & 164~s & 18~s & Default 120~s timeout; full behavioural diff + SHA-256 verify \\
        \bottomrule
    \end{tabularx}
\end{table}

\section{Tests}
\label{sec:testes}

\textbf{Functional behaviour.}
Automated coverage totals \textbf{220 tests}: 110 backend pytest cases (API contracts, scoring,
YARA compilation, pipeline logging, sandbox drivers), 62 frontend Vitest unit tests, Playwright
end-to-end upload flows, and 48 .NET xUnit tests (vm-agent, WPF, benign smoke binary). CI also
verifies YARA rule compilation and diagram synchronisation. Path~B was exercised with a dedicated
benign smoke-test binary; Path~A validates execution gating and exit-code reporting.

\textbf{Design rationale.}
Tests were distributed across layers because failures in malware platforms often appear at
interfaces---upload validation, job state transitions, VM driver stubs---rather than in a single
module. pytest with httpx exercises the API without a browser; Vitest and Playwright cover UI
contracts that backend tests cannot see; xUnit covers .NET components that Python cannot reach.
Synthetic scoring profiles and benign PE corpora (Section~\ref{sec:discussion-results}) complement
unit tests by validating end-to-end scoring behaviour. This pyramid trades some duplication for
confidence that refactors do not silently break the student-facing workflow.

\section{Conclusions}

The implementation realises the architecture planned in Chapter~\ref{ch:engenharia}: a modular static
pipeline, asynchronous jobs, two sandbox orchestration paths, and converging web/desktop entry points.
The dual-path sandbox design reflects a deliberate trade-off---API integration versus instrumentation
depth---rather than accidental duplication. Appendix~\ref{lst:run-gate} illustrates representative
sandbox controls at source level for readers who wish to inspect implementation detail.
